\selectlanguage{ngerman}%
\frenchspacing%
\begin{ledgroupsized}[r]{120mm}%
\footnotesize%
\pstart%
\noindent%
\textbf{Überlieferung:}%
\pend%
\end{ledgroupsized}%
\begin{ledgroupsized}[r]{114mm}%
\footnotesize%
\pstart%
\parindent -6mm%
\makebox[6mm][l]{\textit{L}}%
Notiz:
LH XXXVII\;4 Bl.\;94.
Ein Zettel (5\:x\:6,5\;cm);
Ränder unregelmäßig beschnitten;
Papiererhaltungsmaßnahmen.
Eine Seite auf Bl.\;\href{\Scan{LH_37_04_094r}}{94\,r\textsuperscript{o}}.
Bl.\;\href{\Scan{LH_37_04_094v}}{94\,v\textsuperscript{o}}
ist unbeschrieben.
\pend%
\end{ledgroupsized}%
\vspace{5mm}%
\begin{ledgroup}%
\footnotesize%
\pstart%
\noindent%
\textbf{Datierungsgründe:}
\pagelabel{60249_dat}
Die undatierte Notiz N.\;\ref{60249} ist auf einem Träger überliefert, der keine Spuren eines Wasserzeichens %
aufweist. %
Das Zettelformat deutet aber auf eine Entstehung %
nach Leibnizens Umzug von Paris\indexO{Paris} nach Hannover:\indexO{Hannover}
Kleine und fragmentarische Zettel sind im Nachlass seltener aus der Mainzer\indexO{Mainz} oder Pariser\indexO{Paris} Zeit überliefert.
Demnach ist unwahrscheinlich, dass die Notiz N.\;\ref{60249} vor 1677 verfasst wurde.
\pend%
\pstart%
Weitere Anhaltspunkte, die zum heutigem Zeitpunkt eine begründete Einschränkung der Datierung ermöglichen würden, bietet die Notiz N.\;\ref{60249} sonst keine.
\pend%
\end{ledgroup}%
\selectlanguage{latin}%
\frenchspacing%
\vspace{8mm}%
\normalsize%
\pstart%
\noindent%
\lbrack%
\href{\Scan{LH_37_04_094r}}%
{94\,r\textsuperscript{o}}%
\rbrack\
Cur saepe statim post tonitru%
\indexS{tonitru}
pluvia fit acrior?%
\indexS{pluvia}
\pend%
\pstart%
Post lumen majus tonitru majus;%
\indexS{lumen tonitrus}
videndum
cur tonitru minus post lumen magnum sit remotius,%
\indexS{tonitru}%
\indexS{lumen tonitrus}
seu majoris intervalli.%
\indexS{intervallum inter fulmen et tonitru}
\pend%
\pstart%
\edtext{%
Vulgo ignem ictu fulminis accensum,%
\indexS{ignis accensus fulmine}%
\indexS{ictus fulminis}%
\indexS{fulmen}
non aqua sed lacte extingui debere%
\indexS{aqua ad ignem extinguendum}%
\indexS{lac}
ridicule putant.%
}{%
\lemma{Vulgo \lbrack...\rbrack\ putant}%
\Cfootnote{%
Anspielung nicht erschlossen.%
}}
\pend%
\selectlanguage{latin}%
\frenchspacing%
\normalsize%
\newpage%
